\section{Open Questions (Follow-on Probes)}

\begin{enumerate}
\item \textbf{Quantitative depth-reduction crossover $p^\star$.}
At what QAOA depth $p^\star$ does the 12-qubit binary encoding match
the qutrit encoding's approximation ratio, and does $p^\star/p_{\text{qudit}}$
scale with $N$ or $K$? Our results show qudit strictly beats qubit at
every $p\in\{1..5\}$, but the paper never fits a crossover; without
$p^\star$ the ``depth-reduction'' claim is only directional, not
quantitative. \emph{Next steps:} extend the simulator to $p\in\{6..12\}$
(still $<729{\times}729$, cheap) at $N\in\{4..8\}$ and $K\in\{3,4,5\}$;
fit $p^\star(N,K)$ so binary QAOA at $p^\star$ equals qudit at fixed
$p_q{=}3$; publish the ratio.

\item \textbf{Non-Pauli qudit noise robustness.}
How does qutrit QAOA degrade under realistic non-Pauli qudit noise
(e.g. spontaneous emission between angular-momentum levels, or Kraus
channels that leak out of the $d{=}3$ subspace)? The paper motivates
qudit QAOA for cold-atom / trapped-ion hardware where errors are
\emph{not} Pauli, but our replication is pure state-vector.
\emph{Next steps:} add a Lindblad-master-equation step with a
parametric $d{=}3\to d{=}4$ leakage rate and a per-gate depolarizing
channel; sweep noise strength and re-measure the C3 gap ratio.

\item \textbf{Beyond $k=3$: general max-$k$-coloring / max-$k$-cut.}
Does the qudit advantage survive at $k\in\{4,5,6\}$, or does the
$k=3$ result exploit SU(3)-specific algebra? Only $k=3$ was tested
here and in the paper's Sec.~IV headline; larger $k$ uses ququarts /
qupents whose $X+X^\dagger$ eigenspectrum has different structure.
\emph{Next steps:} reimplement for general $d$; run $k=4,5,6$ coloring
instances at matched depth vs.\ multi-qubit binary encodings; check
whether the gap ratio grows or shrinks with $k$.

\item \textbf{Adaptive-depth qudit QAOA (ADAPT-QAOA-style).}
Can an adaptive-depth qudit variant reach the same ground-state mass
with fewer parameters than the fixed-depth $p{=}5$ baseline used here
and in the paper? Standard QAOA has $2p$ parameters; ADAPT-QAOA grows
the ansatz operator-by-operator by gradient. No qudit ADAPT variant
appears in the paper. \emph{Next steps:} extend the mixer pool from
$\{X+X^\dagger\}$ to $\{X+X^\dagger, Z, XZ, ZX, \dots\}$; greedily
pick the operator with the largest energy gradient; compare parameter
count to reach $P(\text{gs})\ge 0.85$ vs.\ fixed $p{=}5$.

\item \textbf{Hybrid qudit-qubit encodings for real problem features.}
Does a hybrid qudit-qubit encoding (one qutrit for color, extra qubits
for auxiliary constraints such as EV charging-station capacity)
recover the qudit gap advantage while enabling problem features that
pure qudits cannot express? The paper's EV-charging motivation has
integer color variables \emph{and} binary capacity constraints; the
latter is unnatural in pure qudits. \emph{Next steps:} add binary
ancilla qubits with a mixed $d{=}3\times d{=}2\times\dots$ Hilbert
space; encode a capacity constraint via ancilla-conditioned penalty;
compare gap and $P(\text{gs})$ to pure-qutrit and pure-qubit baselines.
\end{enumerate}
